\Needspace{9\baselineskip}
\section{Homogeneous sections on a Hopf manifold}\label{sec:hopf-homogeneous-sections}

On a scalar Hopf manifold, a holomorphic section of a line bundle lifts to a function on $\mathbb C^n\setminus\{0\}$ with an exact homogeneity property under the contraction. Hartogs extension reduces the section spaces to spaces of homogeneous polynomials, and the invariance of the metric makes the Bergman kernel constant. The resulting growth is of order $m^{n-1}$, reflecting the dimension of the projective base rather than that of the manifold.

Let $n\geq2$, let $z=(z_1,\ldots,z_n)$ be the coordinates on $\mathbb C^n$, and set
\[
 X_n=(\C^n\setminus\{0\})/\langle z\mapsto \mathrm{e}^{-1}z\rangle.
\]

\Needspace{6\baselineskip}
On the covering space put
\[
 s=\log|z|^2,\qquad \alpha=\ddc s,\qquad
 \zeta=\mathrm{d}s\wedge\dc s,\qquad \omega=\alpha+\zeta.
\]
The contraction changes $s$ by the constant $-2$, so these forms descend to $X_n$. The calculation proving \eqref{eq:hopf-torsion}, written in \cref{chap:hopf} for $n=3$ and with $\beta$ in place of $\zeta$, applies in every dimension and gives
\[
 \alpha^n=0,\qquad \zeta^2=0,\qquad
 \ddc\zeta=-\alpha^2,
 \qquad \ddc(\omega^{n-1})=-(n-1)\alpha^n=0.
\]
Thus $\omega$ is a Gauduchon form. Projectivization defines a holomorphic map $\pi\colon X_n\to\mathbb P^{n-1}$ with elliptic fibers, and $\alpha$ is the pullback of the Fubini--Study form normalized to have total mass one. Since $\zeta$ restricts to twice the normalized area form on each fiber, as in the proof of \cref{prop:hopf-mass}, and $\omega^n=n\alpha^{n-1}\wedge\zeta$, we obtain
\[
 \int_{X_n}\alpha^{n-1}\wedge\zeta=2,
 \qquad \int_{X_n}\omega^n=2n.
\]
Set $\mathrm{d}V=\omega^n/(2n)$, a probability measure.

For an integer $k\geq0$, let $L_k$ be the line bundle on $X_n$ obtained from $(\C^n\setminus\{0\})\times\C$ through the identification
\[
 (z,t)\sim(\mathrm{e}^{-1}z,\mathrm{e}^{-k}t).
\]
Give it the metric $|t|_{h_k}^2=|t|^2|z|^{-2k}$. Under the identification the two factors change by $\mathrm{e}^{-2k}$ and $\mathrm{e}^{2k}$, so this metric descends. In the frame $t=1$ the squared norm is $|z|^{-2k}$, and the curvature form is $k\alpha$. We apply the constructions of \cref{sec:weighted-section-spaces} with $X=X_n$, $L=L_k$, $h=h_k$ and this probability volume form.

\Needspace{15\baselineskip}
\begin{proposition}\label{prop:hopf-bergman-exact}\index{Bergman kernel!homogeneous Hopf example}
For $m\geq1$ and the weight $u=0$,
\[
 N_m:=h^0(X_n,L_k^m)=\binom{km+n-1}{n-1},
 \qquad B_{m,0}=N_m.
\]
The Bergman potential is the constant $m^{-1}\log N_m$, and the Bergman current is $k\alpha$. Moreover,
\[
 \deg_\omega L_k=2k(n-1),\qquad
 \lim_{m\to\infty}m^{-n}N_m=0,
 \qquad (k\alpha)^n=0,
\]
whereas $\int_{X_n}\omega^n=2n>0$.
\end{proposition}
\begin{proof}
A section of $L_k^m$ lifts to a holomorphic function $f$ on $\C^n\setminus\{0\}$ satisfying
\[
 f(\mathrm{e}^{-1}z)=\mathrm{e}^{-km}f(z).
\]
Since $n\geq2$, Hartogs' theorem extends $f$ holomorphically to $\C^n$. In its Taylor series, a monomial of total degree $j$ is multiplied by $\mathrm{e}^{-j}$ under the contraction, so comparing coefficients shows that only monomials of degree $km$ occur. Conversely, every homogeneous polynomial of that degree has the required equivariance. Counting monomials gives the binomial coefficient; when $k=0$, only the constant polynomials remain.

The inner product defined by $h_k^m$ and $\mathrm{d}V$ is invariant under the unitary group $U(n)$, so the basis-independent function $K_{m,0}$ is $U(n)$-invariant. It is also invariant under positive radial dilations: every section is homogeneous of degree $km$, and the factor $|z|^{-2km}$ in its squared norm cancels the effect of a dilation. Unitary transformations and positive radial dilations act transitively on $\C^n\setminus\{0\}$, so $K_{m,0}$ is constant. By \cref{prop:finite-bergman-identities} its integral is $N_m$, and $\mathrm{d}V$ has mass one. Thus $K_{m,0}=B_{m,0}=N_m$.

\Needspace{6\baselineskip}
The statements about the Bergman potential and current follow immediately. Since $\alpha^n=0$ and $\zeta^2=0$,
\[
 \int_{X_n}k\alpha\wedge\omega^{n-1}
 =k(n-1)\int_{X_n}\alpha^{n-1}\wedge\zeta
 =2k(n-1).
\]
The binomial coefficient grows at most like $m^{n-1}$, which proves the normalized limit, and $\alpha^n=0$ proves the vanishing of the curvature mass.
\end{proof}

For $n=3$, $X_3$ is the Gauduchon threefold of \cref{thm:hopf-unbounded}, on which Hermitian Monge--Amp\`ere masses are unbounded; for $k>0$, the line bundles $L_k$ nevertheless have positive degree and quadratic section growth. In the notation of \cref{ex:hopf-exact-complement-index}, $L_k$ is the line bundle $E_a$ on $X_q$ with $q=\mathrm e^{-1}$ and $a=\mathrm e^{-k}$. Hence, for every $n\geq2$ and $k>0$, the complement $\mathbb P(\mathcal O\oplus L_k)\setminus\mathbb P(\mathcal O)$ of \cref{subsec:flat-restriction-hopf-complements} has least smooth convexity index two, that is, it is $2$-convex but not $1$-convex in the sense of \cref{subsec:concave-neighbourhood-dual-norms}. When $k=0$, the line bundle is trivial and the least smooth convexity index of its complement is $n+1$, by the same example.
